\BourbakiExercises{Exercises}

\begin{enumerate}
\def\labelenumi{\arabic{enumi})}
\tightlist
\item
  Let A and B be K-algebras and M an A-module; view \(M \otimes B\) as a module over the algebra \(C = A \otimes B\) . Identify M with an A-submodule of \(M \otimes B\) .
\end{enumerate}

\begin{enumerate}
\def\labelenumi{\alph{enumi})}
\item
  Prove the inclusion \(M \cap \mathfrak{R}_{\mathrm{C}}(\mathrm{M} \otimes \mathrm{B}) \subset \mathfrak{R}_{\mathrm{A}}(\mathrm{M})\) (consider an A-linear homomorphism from M to a simple A-module).
\item
  Prove that we have equality in each of the following cases:
\end{enumerate}

\begin{enumerate}
\def\labelenumi{(\roman{enumi})}
\item
  A is finite-dimensional over \(\mathbf{K}\) .
\item
  The A-module M is finitely generated, and A is the union of a right directed family of subalgebras that are finite-dimensional over K.
\item
  \(M = A_{s}\) , and the radical of A is nilpotent.
\end{enumerate}

(In case (i), one proves that every homomorphism from \(\mathrm{M} \otimes \mathrm{B}\) to a simple C-module S is zero on \(\Re_{\mathrm{A}}(\mathrm{M})\) , by observing that \(\Re_{\mathrm{C}}(\mathrm{S})\) is zero.)

\begin{enumerate}
\def\labelenumi{\arabic{enumi})}
\setcounter{enumi}{1}
\item
  Let \(\mathrm{A}\) be a K-algebra that is an integral domain and \(\mathrm{L}\) its field of fractions. Prove that \(\mathrm{L}\) is a simple \(\mathrm{A} \otimes \mathrm{L}\) -module but that its radical as an A-module is equal to \(\mathrm{L}\) .
\item
  Give an example of (nonzero) K-algebras \(\mathrm{A}_1\) and \(\mathrm{A}_2\) such that \(\Re (\mathrm{A}_1)\neq 0\) and \(\Re (\mathrm{A}_1\otimes \mathrm{A}_2) = 0\) (cf.~VIII, p.~167, Exercise 11).
\item
  Give an example of a commutative field K and two extensions E and F of K such that the ring \(A = E \otimes_{K} F\) is semisimple but is not a field (resp. is not semisimple). Deduce that the A-module \(A_{s} = E_{s} \otimes_{K} F_{s}\) is semisimple but not simple (resp. is not semisimple).
\end{enumerate}

¶ 5) Let A be a quasi-simple K-algebra (VIII, p.~120, Remark 2).

\begin{enumerate}
\def\labelenumi{\alph{enumi})}
\item
  Prove that the (A, A)-bimodule \(_{s}A_{d}\) is simple; deduce that the center Z of A is a field, isomorphic to the commutant of this bimodule.
\item
  Let B be a K-algebra. Prove that every two-sided ideal of \(A \otimes B\) is of the form \(A \otimes_{Z} b\) , where b is a two-sided ideal of B (reduce to the case Z = K, and use Corollary 2 of VIII, p.~63).
\item
  Deduce that if A and B are quasi-simple K-algebras of which one has center K, then the algebra \(\mathrm{A} \otimes \mathrm{B}\) is quasi-simple.
\end{enumerate}

\(d\) ) Let B be a K-algebra and A a quasi-simple subalgebra of B whose center Z contains K. Prove that A and its commutant \(\mathrm{A}'\) in B are linearly disjoint over Z (use \(b\) ).

\begin{enumerate}
\def\labelenumi{\alph{enumi})}
\setcounter{enumi}{4}
\item
  Let A be a quasi-simple K-algebra and Z its center. Let \(\varphi : \mathrm{A} \otimes \mathrm{A}^{\circ} \to \operatorname{End}_{\mathrm{Z}}(\mathrm{A})\) be the homomorphism that sends \(a \otimes b\) to the endomorphism \(x \mapsto axb\) . Prove that \(\varphi\) is injective and that its image H is a quasi-simple algebra (use \(c\) ) and \(d\) ). Prove, moreover, that H is dense in \(\operatorname{End}_{\mathrm{Z}}(\mathrm{A})\) (VIII, p.~129, Exercise 3); we have \(\mathrm{H} = \operatorname{End}_{\mathrm{Z}}(\mathrm{A})\) if and only if A is finite-dimensional over Z (cf.~VIII, p.~128, Exercise 2).
\item
  Let A be a semisimple algebra with center K. Deduce from e) that the K-algebra \(A \otimes A^{o}\) is semisimple if and only if A has finite rank over K (cf.~VIII, p.~238, Theorem 2).
\end{enumerate}

\begin{enumerate}
\def\labelenumi{\arabic{enumi})}
\setcounter{enumi}{5}
\item
  Let A be a K-algebra with nonzero nilpotent radical. Prove that the homomorphism \(\varphi: A \otimes A^{\circ} \to \operatorname{End}_{K}(A)\) (Exercise 5) is not injective.
\item
  Let \(A_{1}\) and \(A_{2}\) be semisimple K-algebras and \(Z_{1}\) and \(Z_{2}\) their centers. Suppose that \(A_{1}\) has finite degree over K and that the ring \(Z_{1} \otimes Z_{2}\) is reduced. Let f be a K-algebra homomorphism from \(A_{2}\) to \(A_{1}\) . Prove that the commutant of \(f(\mathrm{A}_{2})\) in \(A_{1}\) is a semisimple K-algebra (view \(A_{1}\) as an \((\mathrm{A}_{2} \otimes \mathrm{A}_{1}^{\circ})\) -module and use Proposition 7 of VIII, p.~221).
\item
  Let K be a commutative field and L a radical extension of K of finite degree \textgreater{} 1. Let V be an infinite-dimensional vector space over L, and let A be the K-subalgebra of \(\operatorname{End}_{\mathrm{L}}(\mathrm{V})\) generated by the endomorphisms of finite rank and the identity mapping. Prove that the L-algebra \(\mathrm{A}_{(\mathrm{L})}\) has nonzero radical even though its center is equal to L (use Exercise 7).
\item
  Let A and B be K-algebras, M a free A-module, and N a B-module.
\end{enumerate}

\begin{enumerate}
\def\labelenumi{\alph{enumi})}
\item
  Prove that the image of the canonical homomorphism \(\vartheta\) from \(\mathrm{End}_{\mathrm{A}}(\mathrm{M})\otimes \mathrm{End}_{\mathrm{B}}(\mathrm{N})\) to \(\mathrm{End}_{\mathrm{A}\otimes \mathrm{B}}(\mathrm{M}\otimes \mathrm{N})\) is a dense subring of the ring \(\mathrm{End}_{\mathrm{A}\otimes \mathrm{B}}(\mathrm{M}\otimes \mathrm{N})\) .
\item
  If the ring of homotheties of the B-module N is dense in its bicommutant, then the same holds for the ring of homotheties of the \((A \otimes B)\) -module \(M \otimes N\) .
\item
  Suppose that there exist an element y of N and a sequence \((v_{n})\) of endomorphisms of N such that the elements \(v_{n}(y)\) generate an infinite-dimensional vector space over K. Prove that the mapping \(\vartheta\) is not surjective.
\end{enumerate}

\begin{enumerate}
\def\labelenumi{\arabic{enumi})}
\setcounter{enumi}{9}
\tightlist
\item
  Keep the assumptions of Exercise 9 and suppose, moreover, that the B-module N is simple; denote its commutant by D.
\end{enumerate}

\begin{enumerate}
\def\labelenumi{\alph{enumi})}
\item
  The mapping \(\vartheta\) is surjective if and only if M admits a finite basis over A or D is finite-dimensional over K (use Exercise 9, c)).
\item
  Prove that the \((\mathrm{A} \otimes \mathrm{B})\) -module \(M \otimes N\) is semisimple if and only if the ring \(A \otimes D^{o}\) is semisimple (reduce to the case \(M = A_{s}\) ; observe that \(M \otimes N\) is a free \(A \otimes D^{o}\) -module, and use Proposition 5 of VIII, p.~85 as well as Exercise 9, a)).
\end{enumerate}

¶ 11) Let A and B be K-algebras; suppose that the rings A and B are primitive and that their respective socles s and t (VIII, p.~130, Exercise 8) are not reduced to 0. Let D (resp. E) be the commutant of a minimal left ideal a of A (resp. b of B); suppose that the algebra D ⊗ E is simple. Prove that A ⊗ B is a primitive ring with socle s ⊗ t (observe that the (A ⊗ B)-module a ⊗ b is isotypical of finite length; use Exercise 9, a) of VIII, p.~131).

\begin{enumerate}
\def\labelenumi{\arabic{enumi})}
\setcounter{enumi}{11}
\tightlist
\item
  Let \(A_{1}\) be the K-algebra \(\mathrm{K}(\mathrm{X})[\mathrm{U}]_{1,d/d\mathrm{X}}\) (VIII, p.~11), \(A_{2}\) the K-algebra \(K[[T]]\) , and A the K-algebra \(A_{1} \otimes_{K} A_{2}\) . Identify \(A_{1}\) and \(A_{2}\) with subalgebras of A.
\end{enumerate}

\begin{enumerate}
\def\labelenumi{\alph{enumi})}
\item
  Prove that \(A_{2}\) is the center of A, that A does not contain any zero divisors, and that the invertible elements of A are those of \(A_{2}\) .
\item
  Prove that the two-sided ideals of A are the ideals \(AT^{n}\) for \(n \geqslant 0\) .
\item
  Prove that the ring A is primitive even though its center has nonzero radical (observe that a maximal left ideal of A containing 1 - TU does not contain any two-sided ideals other than 0).
\end{enumerate}

\begin{enumerate}
\def\labelenumi{\arabic{enumi})}
\setcounter{enumi}{12}
\tightlist
\item
  Let D and E be fields with centers containing K, and let V (resp. W) be a vector space over D (resp. E). Prove that the \((\mathrm{D} \otimes_{\mathrm{K}} \mathrm{E})\) -module \(V \otimes_{K} W\) is free.
\end{enumerate}

¶ 14) Let V be a vector space over a field D. Denote the endomorphism ring of the abelian group V by \(\Omega\) , and identify D with a subring of \(\Omega\) . Set \(A = \text{End}_{\text{D}}(\text{V})\) .

\begin{enumerate}
\def\labelenumi{\alph{enumi})}
\item
  Prove that A and D are linearly disjoint over their common center Z (cf.~Exercise 5, d)).
\item
  The image of \(D \otimes_{Z} A\) in \(\Omega\) is equal to \(\operatorname{End}_{\mathbb{Z}}(\mathrm{V})\) if and only if D has finite rank over Z. (To see that the condition is necessary, first deduce the case V = D from Exercise 5, e), and then fix an element \(x \neq 0\) of V and consider the subalgebra B of \(\operatorname{End}_{\mathbb{Z}}(\mathrm{V})\) consisting of the Z-endomorphisms u of V such that \(u(\mathrm{D}x) \subset \mathrm{D}x\) . Observe that the equality \(\operatorname{End}_{\mathbb{Z}}(\mathrm{V}) = \mathrm{DA}\) implies \(\mathrm{B} = \mathrm{D}(\mathrm{B} \cap \mathrm{A})\) .)
\end{enumerate}

\begin{enumerate}
\def\labelenumi{\arabic{enumi})}
\setcounter{enumi}{14}
\tightlist
\item
  \begin{enumerate}
  \def\labelenumii{\alph{enumii})}
  \tightlist
  \item
    Let \(\mathrm{K}\) be a commutative field, A a K-algebra, and \(\mathrm{B} = \mathrm{K}[\mathrm{X}_{\iota}]_{\iota \in \mathrm{I}}\) a polynomial algebra over \(\mathrm{K}\) . Prove that if every nonzero element of \(\mathrm{A}\) is regular, then the algebra \(\mathrm{A} \otimes_{\mathrm{K}} \mathrm{B}\) has the same property.
  \end{enumerate}
\end{enumerate}

\begin{enumerate}
\def\labelenumi{\alph{enumi})}
\setcounter{enumi}{1}
\tightlist
\item
  Deduce that if D is a field of finite rank over K and E is a purely transcendental extension of K, then the ring \(D \otimes_{K} E\) is a field.
\end{enumerate}

¶ 16) Let D be a field, V an infinite-dimensional right vector space over D, and A a dense subring of the ring \(\mathrm{End}_{\mathrm{D}}(\mathrm{V})\) whose socle \(\mathfrak{s}\) is not zero (VIII, p.~130, Exercises 8 and 9). Let K be a subfield of the center of D and B a K-algebra. Prove that if the radical of the K-algebra \(\mathrm{D}^{\circ} \otimes \mathrm{B}\) is nonzero, then so is that of \(\mathrm{A} \otimes \mathrm{B}\) . (Let \(\mathfrak{a}\) be the set of elements \(\sum u_{i} \otimes b_{i}\) of \(\mathfrak{s} \otimes \mathrm{B}\) satisfying \(\sum \langle u_{i}(x), x^{*} \rangle \otimes b_{i} \in \Re(\mathrm{D}^{\circ} \otimes \mathrm{B})\) for all \(x \in \mathrm{V}\) and \(x^{*} \in \mathrm{V}^{*}\) . Prove that \(\mathfrak{a}\) is a two-sided ideal of \(\mathrm{A} \otimes \mathrm{B}\) and, using Exercises 10, e) of VIII, p.~167 and 9 of VIII, p.~131, that it is nonzero and is contained in the radical of \(\mathrm{A} \otimes \mathrm{B}\) .)

\begin{enumerate}
\def\labelenumi{\arabic{enumi})}
\setcounter{enumi}{16}
\tightlist
\item
  Let K be a commutative field and D and E fields whose centers contain K. Suppose that \([E:K]=m\) is finite and that the center of one of the fields D and E is equal to K. The K-algebra \(D\otimes_{K}E\) is then isomorphic to the endomorphism algebra of a finite-dimensional vector space h over a field C (VIII, p.~222, Corollary 2).
\end{enumerate}

\begin{enumerate}
\def\labelenumi{\alph{enumi})}
\item
  Prove that \(h\) divides \(m\) (cf.~VIII, p.~205, formula (32)).
\item
  Let W be a vector space over D; then E is isomorphic to a K-subalgebra of \(\operatorname{End}_{\mathrm{D}}(\mathrm{W})\) if and only if the dimension of W is infinite or a multiple of m/h (note that W is then a \((\mathrm{D} \otimes_{\mathrm{K}} \mathrm{E})\) -module).
\item
  Prove that if \([\mathrm{D}:\mathrm{K}]\) is finite and prime to \(m\) , then \(\mathrm{D}\otimes_{\mathrm{K}}\mathrm{E}\) is a field (apply \(a\) ).
\end{enumerate}

